\begin{abstract}
\porozhdeno{\Claim{A01}{Сравнены готовые модели, связанные с базами SILVA, GTDB и Greengenes2, при одном алгоритме таксономического назначения и фиксированной обработке ампликонов почв Москвы. Для \ProjectSamplesSoil{} почвенных образцов оценены покрытие назначениями, буквальное совпадение родословных и расстояние профилей. Медиана расстояния семейственных профилей SILVA--GTDB равна \CmpSilvaGtdbFamilyDistanceMedian{} при межобразцовом диапазоне \ArticleSilvaGtdbFamilyDistanceMin{}--\ArticleSilvaGtdbFamilyDistanceMax{} и медианах покрытия ридов \CmpSilvaFamilyReadMedianPct{}\% и \CmpGtdbFamilyReadMedianPct{}\% соответственно; неназначенная масса включена в профили. Отдельный отбор бактерий и архей меняет знаменатели и может оставлять пустой набор либо неназначенное семейство при непустом наборе. Числа проверены по сохранённым результатам независимым скриптом. Различия показывают чувствительность назначений к готовой модели, но без эталона не устанавливают точность и не объясняют состав почв экологическими причинами.}}
\end{abstract}
